\chapter{Parser}
\label{ch:parser}

\section{Overview}

The parser (\texttt{zparse.c}) performs recursive-descendant parsing over the token stream produced by the lexer and builds an AST (Abstract Syntax Tree) composed of \texttt{ZNode} nodes.

\section{ZNode}
Every construct in the language is represented as a tagged union node:
\begin{verbatim}
struct ZNode {
    ZNodeType type;
    union {
        /* type-specific fields */
    };
}
\end{verbatim}

\section{Parsing Strategy}
The parser is hand-written recursive descent parser. Operator precedence is handled via the standard technique of one parsing function per precedence level, from lowest (\texttt{parseExpr}) to highest (\texttt{parsePrimary}).
The parser also has some helpers to handle the token stream:
\begin{itemize}
    \item{\texttt{peek}} returns the current token
    \item{\texttt{consume}} returns the current token and consumes it
    \item{\texttt{check}} returns true if the current token is the same type of the passed type
    \item{\texttt{match}} same as \texttt{check} but if true consumes the token
    \item{\texttt{expect}} invalidate the function (by returning \texttt{NULL}) if the token doesn't match
\end{itemize}

Before parsing the entire token stream, the parser search macro definitions and stores it in a list inside \texttt{ZParser}.
The parser starts parsing file-level statements calling one of the following functions depending on the current token.
\begin{itemize}
    \item{\texttt{parseFuncDecl}} parse a function declaration
    \item{\texttt{parseFuncBlock}} parse a block of function (static block or receiver block)
    \item{\texttt{parseEnumDecl}} parse an enum declaration
    \item{\texttt{parseStructDecl}} parse a struct declaration
    \item{\texttt{skipMacro}} skip the token of a macro saved in the first phase
\end{itemize}

Inside a block, \texttt{parseBlock} dispatches to statement-level parser including \texttt{parseIf,parseLoops,parseMatch,parseReturn,parseDefer,parseVarDef}.
\section{Inline blocks (\zn{->})}
A function body may be written as a signel inline expression with \texttt{->}:
\zn{square :: i32 (i32 x) -> x * x}

\texttt{parseBlockOrInline} handles this by detecting the \zn{->} token and wrapping the expression in a \texttt{NODE\_RETURN} node so the semantic and codegen passes treat it uniformly.

Also \zn{if} and \zn{for} supports the inline block but they are converted in a normal \texttt{NODE\_BLOCK} with a single statement (the parsed expression).

\begin{note}
    \zn{if} can be used as expression via the \zn{->} operator. For more details chek the \ref{sec:flow}{language reference}.
\end{note}

\section{Error Recovery}
The parser is the most fagile phase in the pipeline: a single malformed statement acn desynchronise the token stream and trigger a cascade of spurious errors.

The parser also needs a limited form of backtracking for positions that are syntactically ambiguous. When the parser encounters \zn{if}, for instance, it cannot yet determine whether it is parsing an if statement or an if expression. The strategy is to attempt each alternative in order and return the result of the first one that succeeds. This is the role of \texttt{parseOrGrammar}: before trying an alternative, it saves the current stream position (calling \texttt{store}), if that altenative fails, it restores the position (calling \texttt{undo}) and tries the next one.
Errors are managed through a separate bookkeeping structure. The parser maintains a stack of \texttt{breakpoints}: when entering an alternative, a breakpoint is pushed; if parsing succeeds the breakpoint is committed, otherwise it is popped and the accumulated errors for that attempt are discarded.

The most difficult aspect of error recovery is handling an invalid statement, The conventional approahc is to skip tokens until a known terminato is found\} for blocks, \zn{;} for statements and then resume parsing from that point. In Zinc, however, statements have no explicit terminator, which makes it significantly harder to locate a safe resynchronisation point.
